\section{把诉求写成目标函数与约束}
\label{sec:18-Synthesis/01-System-Lifecycle/01}

会议室里，运营负责人说的是一句话：``我们想提升用户满意度。''

散会之后，这句话要变成一行代码里的 \texttt{loss}。中间发生了什么，会议纪要上一个字都没有。有人想起满意度调研的回收率只有百分之三，于是改用``七日内是否再次下单''当作代理；有人指出复购受价格活动影响太大，于是加上``会话结束后是否点了好评''；两个信号权重怎么定？没人有依据，于是各取一半。客服团队顺口提了一句``别让平均处理时长涨太多''，这句话被写成一个软性罚项，系数拍了个 \(0.1\)。

三个月后系统上线，平均处理时长确实没涨，好评率涨了四个百分点，复购涨了一个多点。同一时期客诉量涨了三成。复盘时才发现，模型学会了在会话末尾插入一句诱导好评的话术，并且倾向于把复杂问题快速转给人工——因为``转人工''会让这段会话提前结束，处理时长的账记不到它头上。

这一步是整条生命周期里最难的，也是唯一一步没有任何算法可以代劳的。它难不是因为技术，是因为\textbf{要把一个模糊的愿望压缩成一个可以被机械最大化的数，压缩过程中丢掉的东西不会消失，它们会在下游以故障的形态回来}。

\subsection{代理指标与真实目标之间的缝隙会被优化压力撑开}

真实目标——用户是不是真的满意——不可测量。可测量的是它的代理：好评率、复购率、会话时长。代理之所以能当代理，是因为在\emph{当前的行为分布下}它与真实目标高度相关。但优化恰恰是要改变行为分布的。一旦系统开始朝代理指标施加压力，那个相关性赖以成立的机制就被破坏了。

\begin{center}
\begin{tikzpicture}[x=1cm,y=1cm,font=\small,>=stealth]
  \draw[->,black!60] (0,0) -- (11.6,0) node[anchor=north east,font=\footnotesize,black!70] {施加在代理指标上的优化压力};
  \draw[->,black!60] (0,0) -- (0,4.3);
  \node[rotate=90,anchor=south,font=\footnotesize,black!70] at (-0.45,2.1) {指标水平};
  % 代理指标：单调上升
  \draw[very thick,blue!65!black,smooth] plot coordinates
    {(0,0.9) (1.5,1.75) (3,2.5) (4.5,3.1) (6,3.55) (7.5,3.85) (9,4.0) (10.6,4.08)};
  % 真实目标：先升后降
  \draw[very thick,orange!80!black,smooth] plot coordinates
    {(0,0.9) (1.5,1.7) (3,2.3) (4.2,2.5) (5.2,2.42) (6.5,2.05) (8,1.5) (9.5,0.95) (10.6,0.65)};
  \draw[dashed,black!45] (4.2,0) -- (4.2,2.5);
  \node[font=\footnotesize,black!70,anchor=north] at (4.2,-0.08) {分道点};
  \node[anchor=west,blue!65!black,font=\footnotesize] at (8.4,4.25) {代理指标（好评率）};
  \node[anchor=west,orange!80!black,font=\footnotesize] at (8.0,0.45) {真实目标（满意度）};
  \draw[<->,black!55] (8,1.55) -- (8,3.85);
  \node[anchor=west,font=\footnotesize,black!70,align=left] at (8.1,2.7) {缝隙};
  \draw[black!45,fill=green!8] (0.35,3.05) rectangle (3.95,4.0);
  \node[align=center,font=\footnotesize,black!75] at (2.15,3.52) {这一段里代理是\\可信的};
\end{tikzpicture}

\emph{代理指标与真实目标在低压力区几乎重合，正是这段重合让人相信代理可用。压力越过分道点之后两条线分开，而监控盘上只画得出蓝线。}
\end{center}

分道点的存在是结构性的，不是执行不力。只优化点击总数会催生夸张标题，只考核平均处理时长会催生提前结单，只看好评率会催生诱导话术——这些不是模型``学坏了''，是模型准确地做了被要求的事。真正被要求的是那个代理，不是那个愿望。

这条经验规律通常被称作 Goodhart 现象，它不是定理，没有精确的成立条件，但它的机制可以说清：代理 \(S\) 与目标 \(U\) 的关系可以写成 \(U = S - G\)，其中 \(G\) 是二者的缺口。在自然波动下 \(G\) 是小的随机项，于是最大化 \(S\) 近似等于最大化 \(U\)；而优化会主动搜索使 \(S\) 大的区域，那正是 \(G\) 大的区域。选择性搜索把一个原本无害的残差项变成了系统性偏差。所以能否用代理，取决于打算施加多大压力，而不只是取决于历史相关系数有多高。

上图那条蓝线还有一层含义：它是唯一被监控的那条。真实目标不可测量意味着它也不可监控，分道之后没有任何仪表会亮红灯。这是下一小节要处理的问题。

\subsection{约束不是限制，是把``什么样的解不可接受''写下来}

上面那个案例里，客服团队那句``别让处理时长涨太多''本来是一条约束，却被写成了系数 \(0.1\) 的罚项。这两种写法不等价，差别在于优化器面对它们的态度。

硬约束 \(g(x)\le 0\) 说的是：违反它的解根本不在考虑范围内，无论目标值多高。罚项 \(f(x)+\lambda g(x)\) 说的是：违反是可以的，只要收益超过 \(\lambda\) 倍的违反量。把安全边界写成罚项，等于给优化器标了一个价格，然后它会理性地按这个价格买走你的边界。\hyperref[sec:09-Convex-Optimization/01-Optimization-Foundations/01]{``把现实选择翻译成优化问题''}把这一步拆成决策变量、目标、可行域三件事分开确定，其中可行域是最容易被省略的一件——因为没有人会主动去写下``什么显然不行''，而优化器恰恰只知道被写下来的东西。

罚项与约束之间的换算不是随意的。\hyperref[sec:09-Convex-Optimization/05-Duality/04]{``对偶变量是约束的影子价格''}给出的判断是：在最优解处，约束对应的对偶变量等于放松这条约束一个单位能带来的目标改善。这句话把``\(\lambda\) 该取多少''从玄学变成了可回答的问题——如果处理时长每多一秒对业务的真实代价是 \(c\)，那么 \(\lambda\) 应当就是 \(c\)，而不是拍出来的 \(0.1\)。反过来，如果解出来的对偶变量远大于业务能承受的那个数，说明这条约束卡得太紧，值得回去和业务方重新谈。多目标加权也是同一件事：权重就是各目标之间的兑换率，把它写下来意味着承认``一个百分点的好评率值多少秒处理时长''，而这个兑换率无论如何都会存在，区别只在于是被明确写下来还是被默认地藏在代码里。

约束还承担另一个角色。可行域的形状决定了这个问题解起来是容易还是绝望：\hyperref[ch:074]{``识别凸优化问题''}的实用价值不在于凸性本身有多美，而在于凸可行域上局部最优即全局最优，于是``解出来了''这句话才有意义。同一个业务约束，写成线性不等式和写成一个非凸集合，后续每一步的可信程度都不一样。这件事应当在建模阶段决定，而不是在调参阶段发现。

\subsection{损失、风险、目标、指标是四个不同的量}

工程语言里这四个词经常混用，混用的后果在评审会上才显现。\hyperref[sec:13-Machine-Learning/01-Learning-Problems-and-Evaluation/02]{``损失、风险、目标与指标不是同一个量''}把它们分开：损失是单个样本上的代价函数，风险是损失在真实分布下的期望，训练目标是风险在有限样本上的经验替代加上正则项，而业务指标是那个真正被汇报的数。

四者两两之间都有缝隙，而且缝隙的性质不同。损失到风险的缝隙是统计问题，靠样本量和评估设计处理；风险到训练目标的缝隙是优化与正则的问题；训练目标到业务指标的缝隙，就是本节开头那个缝隙，它既不是统计问题也不是优化问题，靠加数据和调参永远修不好。把这三种缝隙记在同一本账上，是无数次``模型明明更好了业务却没动''的根源。

再往前一层，损失函数本身就是一个决策论对象。\hyperref[sec:11-Mathematical-Statistics/07-Statistical-Decision-and-Reliable-Prediction/01]{``状态、行动、损失与风险构成一个决策问题''}给出的框架是：先说清楚世界可能处于哪些状态、我们可以采取哪些行动、每个（状态，行动）组合要付出多少，损失函数是这张代价表的写法，不是先验地存在的东西。用平方损失还是绝对损失不是审美偏好——平方损失说的是``偏差两倍时代价四倍''，绝对损失说的是``代价与偏差成正比''，这两句话在风控场景和推荐场景里的真假完全不同。误判成本不对称时（漏检一个欺诈订单和错拦一个正常订单的代价相差两个数量级），却用了对称的损失，那么无论模型多好，它优化的都是别人的问题。

代价表写对之后还有一件事：模型给出的概率要能真的当概率用。\hyperref[sec:11-Mathematical-Statistics/07-Statistical-Decision-and-Reliable-Prediction/06]{``拒绝预测与容量约束把概率变成行动''}处理的是从分数到行动的最后一跳——阈值定在哪里由代价比决定，而人工复核的容量本身就是一条硬约束，它会把``按概率排序''变成一个受限的分配问题。这一步在设计阶段被跳过的话，上线后就会以``模型分数很好但审核团队根本处理不完''的形式回来。

\subsection{这一步做错时，故障长什么样}

第一种形态是指标全线上涨而用户体验下降。所有被监控的数都在变好，客诉、退订、卸载这些没有进目标函数的量在变坏。这几乎总是代理指标与真实目标分道的症状，而且发现得越晚越难回退，因为期间积累的数据已经是被新策略塑造过的。

第二种形态是模型在离线评测上完美，业务方看了一眼说``这不是我要的''。这说明翻译在第一步就偏了：模型解的是一个被精确定义的问题，只不过那个问题不是业务方脑子里的那个。这种故障的特征是双方都拿得出证据，谁也说服不了谁——因为争的其实是那个没被写下来的目标。

第三种形态最隐蔽：约束被悄悄突破而没人注意，因为它当初被写成了罚项。系统在某个子群体上表现异常差，或者在极端输入上给出荒谬输出，回头查会发现从来没有任何一行代码禁止过这件事。

这三种故障都发生在训练之外，也都无法通过更好的模型来修复。它们的共同根源是：目标函数写完之后，所有下游环节都只对它负责，不再对那个原始诉求负责。

诉求变成目标函数之后，下一件要落定的事是数据——用哪批数据、怎样表示。数据看上去比目标客观得多，实际上它携带的假设更多，只是那些假设藏得更深。
